\chapter{Structure de l'Épreuve Écrite et Méthodologie de Réussite}

\section{Format et Modalités de l'Épreuve Écrite}
Les concours de recrutement des Ingénieurs d'État (1\textsuperscript{er} grade) au sein du Ministère de l'Économie et des Finances se caractérisent généralement par une épreuve écrite d'une durée de \textbf{3 heures} (coefficient 4 ou selon l'avis de concours), suivie pour les admissibles d'un entretien oral (coefficient 3).

L'épreuve écrite se présente selon l'une des deux configurations majeures adoptées par le ministère ces dernières années :
\begin{enumerate}
    \item \textbf{Configuration A (Format Prédominant - Sujet de Réflexion / Dissertation Générale \& Économique) :}
    Deux sujets d'actualité stratégique, économique, numérique ou financière sont proposés au choix du candidat, à rédiger en français ou en arabe (ex. session du 28 avril 2024 : Sujet 1 sur la digitalisation / \textit{Maroc Digital 2030}, Sujet 2 sur la transition et sécurité énergétique).
    \item \textbf{Configuration B (Format Hybride - Dissertation + QCM / Questions à Réponses Courtes) :}
    Une partie de culture générale financière, administrative et institutionnelle sous forme de QCM (missions du MEF, Loi Organique de Finances 130-13, finances publiques, fiscalité, déontologie administrative), complétée par une question rédactionnelle d'analyse ou de synthèse technique.
\end{enumerate}

\section{Grille de Notation et Attentes Réelles des Correcteurs du MEF}
Les correcteurs sont des hauts fonctionnaires, directeurs ou inspecteurs du MEF, ainsi que des universitaires. Leurs critères d'appréciation s'articulent autour de cinq axes majeurs :
\begin{table}[H]
\centering
\small
\begin{tabular}{|p{4.5cm}|c|p{8.5cm}|}
\hline
\rowcolor{gray!20} \textbf{Critère d'évaluation} & \textbf{Pondération} & \textbf{Ce que recherche le jury} \\ \hline
\textbf{Structure et Logique du Plan} & 25\% & Problématique claire, plan en deux parties équilibrées (I et II) avec sous-parties (A et B), transitions soignées, respect de la démarche déductive. \\ \hline
\textbf{Maîtrise du Contexte Marocain} & 30\% & Connaissance précise des stratégies nationales (\textit{Maroc Digital 2030}, Charte de l'Investissement, LOF 130-13), citations de lois, chiffres récents, références institutionnelles. \\ \hline
\textbf{Vision de l'Ingénieur d'État} & 20\% & Capacité à lier les enjeux technologiques/techniques à l'efficacité économique, à la rationalisation des coûts et à l'impact citoyen. \\ \hline
\textbf{Qualité Rédactionnelle et Clarté} & 15\% & Style administratif rigoureux, vocabulaire économique précis, orthographe et syntaxe irréprochables, présentation aérée. \\ \hline
\textbf{Esprit Critique et Propositions} & 10\% & Diagnostic lucide des défis et formulation de recommandations réalistes et constructives pour l'administration. \\ \hline
\end{tabular}
\caption{Grille type d'évaluation des épreuves écrites du MEF}
\end{table}

\section{Gestion Optimale du Temps sur 3 Heures (180 Minutes)}
La principale cause d'échec n'est pas le manque de connaissances, mais le manque de temps à la fin de l'épreuve. Voici le découpage chronométré à appliquer rigoureusement :
\begin{itemize}
    \item \textbf{00h00 - 00h15 (15 min) : Choix du sujet et Définition des termes clés}
    Lire attentivement les deux sujets. Choisir celui où l'on dispose du plus grand nombre d'exemples et de chiffres marocains précis. Définir chaque mot du sujet pour éviter le hors-sujet.
    \item \textbf{00h15 - 00h45 (30 min) : Brainstorming, Problématique et Élaboration du Plan détaillé}
    Jeter les idées au brouillon, les classer, formuler la problématique centrale sous forme de question directrice, et construire le plan :
    \begin{center}
    \textbf{I. Diagnostic, État des lieux et Enjeux} (A. Réalisations / B. Limites \& Défis) \\
    \textbf{II. Orientations Stratégiques, Leviers Opérationnels et Perspectives} (A. Cadre stratégique / B. Recommandations)
    \end{center}
    \item \textbf{00h45 - 01h00 (15 min) : Rédaction intégrale de l'Introduction et de la Conclusion au brouillon}
    C'est la première et la dernière impression laissée au correcteur. Elles doivent être parfaites.
    \item \textbf{01h00 - 02h45 (1h45) : Rédaction directe sur la copie d'examen}
    Rédiger directement le corps du texte (Parties I et II) en s'appuyant sur le plan détaillé du brouillon. Soigner les chapeaux introductifs et les phrases de transition.
    \item \textbf{02h45 - 03h00 (15 min) : Relecture impérative}
    Chasse aux fautes d'accord, de ponctuation, vérification de la lisibilité des titres et des chiffres.
\end{itemize}

\section{La Méthode Infaillible de la Dissertation Administrative}

\subsection{L'Introduction en 5 Étapes Indispensables}
Une introduction réussie doit comporter obligatoirement les 5 étapes suivantes en un ou deux paragraphes cohérents :
\begin{enumerate}
    \item \textbf{L'Accroche (Amorce) :} Partir d'une dynamique générale contemporaine, d'un discours royal, d'un événement économique majeur ou d'une mutation globale.
    \item \textbf{Définition et Cadrage du Sujet :} Définir les concepts centraux (ex. digitalisation, transition énergétique, souveraineté économique).
    \item \textbf{Contexte Marocain :} Présenter la place de cette thématique dans les politiques publiques nationales (réformes du MEF, Nouveau Modèle de Développement, Lois de Finances).
    \item \textbf{La Problématique :} Poser la question de fond qui structure toute la réflexion.
    \item \textbf{L'Annonce du Plan :} Annoncer explicitement et élégamment les deux parties principales.
\end{enumerate}

\subsection{La Règle d'Or du Plan en Deux Parties}
La tradition des concours de la fonction publique marocaine et des grands corps d'ingénieurs privilégie le plan binaire équilibré :
\begin{itemize}
    \item \textbf{Partie I : État des Lieux, Acquis et Diagnostic Critique}
    \begin{itemize}
        \item \textit{Sous-partie I.A :} Les réalisations majeures, le cadre institutionnel et les progrès enregistrés au Maroc.
        \item \textit{Sous-partie I.B :} Les contraintes structurelles, les insuffisances et les défis persistants.
    \end{itemize}
    \item \textbf{Partie II : Stratégies Nationales, Rôle de l'Ingénieur et Leviers d'Avenir}
    \begin{itemize}
        \item \textit{Sous-partie II.A :} Les programmes et feuilles de route gouvernementales (ex. stratégies sectorielles, lois-cadres).
        \item \textit{Sous-partie II.B :} Recommandations opérationnelles, modernisation des processus, gouvernance et perspectives d'impact.
    \end{itemize}
\end{itemize}

\subsection{La Conclusion Efficace}
La conclusion doit comprendre deux volets :
\begin{enumerate}
    \item \textbf{Le Bilan (Synthèse) :} Répondre synthétiquement à la problématique sans répéter le texte.
    \item \textbf{L'Ouverture :} Élargir sur un enjeu complémentaire (ex. intégration régionale africaine, souveraineté industrielle, impact de l'intelligence artificielle générative).
\end{enumerate}
